% ============================================================
% Paper Agent — LaTeX 宏包与命令定义
% ============================================================
% 此文件由 main.tex 通过 % ============================================================
% Paper Agent — LaTeX 宏包与命令定义
% ============================================================
% 此文件由 main.tex 通过 % ============================================================
% Paper Agent — LaTeX 宏包与命令定义
% ============================================================
% 此文件由 main.tex 通过 % ============================================================
% Paper Agent — LaTeX 宏包与命令定义
% ============================================================
% 此文件由 main.tex 通过 \input{preamble} 引入。
% 所有 \usepackage、自定义命令、数学符号宏均在此定义。
% Agent 在 glossary.yaml 中新增符号时，必须同步更新此文件。

% --- 中文支持 ---
\usepackage[UTF8]{ctex}

% --- 数学 ---
\usepackage{amsmath}
\usepackage{amssymb}
\usepackage{amsthm}

% --- 图表 ---
\usepackage{graphicx}
\usepackage{booktabs}
\usepackage{caption}
\usepackage{subcaption}

% --- 引用与链接 ---
\usepackage{hyperref}
\usepackage{cleveref}
\usepackage{natbib}

% --- 其他 ---
\usepackage{algorithm}
\usepackage{algorithmic}
\usepackage{xcolor}

% ============================================================
% 自定义数学符号宏
% ============================================================
% 与 config/glossary.yaml 中 symbols 部分保持同步。
% 添加新符号时请同步更新 glossary.yaml。

% \newcommand{\loss}{\mathcal{L}}       % 损失函数
% \newcommand{\params}{\theta}           % 模型参数
% \newcommand{\dataset}{\mathcal{D}}     % 数据集

% ============================================================
% 标记宏 — 用于 Agent 与用户的反馈循环
% ============================================================
\newcommand{\todo}[1]{\textcolor{red}{\textbf{[TODO: #1]}}}
\newcommand{\confirm}[1]{\textcolor{orange}{\textbf{[CONFIRM: #1]}}}
 引入。
% 所有 \usepackage、自定义命令、数学符号宏均在此定义。
% Agent 在 glossary.yaml 中新增符号时，必须同步更新此文件。

% --- 中文支持 ---
\usepackage[UTF8]{ctex}

% --- 数学 ---
\usepackage{amsmath}
\usepackage{amssymb}
\usepackage{amsthm}

% --- 图表 ---
\usepackage{graphicx}
\usepackage{booktabs}
\usepackage{caption}
\usepackage{subcaption}

% --- 引用与链接 ---
\usepackage{hyperref}
\usepackage{cleveref}
\usepackage{natbib}

% --- 其他 ---
\usepackage{algorithm}
\usepackage{algorithmic}
\usepackage{xcolor}

% ============================================================
% 自定义数学符号宏
% ============================================================
% 与 config/glossary.yaml 中 symbols 部分保持同步。
% 添加新符号时请同步更新 glossary.yaml。

% \newcommand{\loss}{\mathcal{L}}       % 损失函数
% \newcommand{\params}{\theta}           % 模型参数
% \newcommand{\dataset}{\mathcal{D}}     % 数据集

% ============================================================
% 标记宏 — 用于 Agent 与用户的反馈循环
% ============================================================
\newcommand{\todo}[1]{\textcolor{red}{\textbf{[TODO: #1]}}}
\newcommand{\confirm}[1]{\textcolor{orange}{\textbf{[CONFIRM: #1]}}}
 引入。
% 所有 \usepackage、自定义命令、数学符号宏均在此定义。
% Agent 在 glossary.yaml 中新增符号时，必须同步更新此文件。

% --- 中文支持 ---
\usepackage[UTF8]{ctex}

% --- 数学 ---
\usepackage{amsmath}
\usepackage{amssymb}
\usepackage{amsthm}

% --- 图表 ---
\usepackage{graphicx}
\usepackage{booktabs}
\usepackage{caption}
\usepackage{subcaption}

% --- 引用与链接 ---
\usepackage{hyperref}
\usepackage{cleveref}
\usepackage{natbib}

% --- 其他 ---
\usepackage{algorithm}
\usepackage{algorithmic}
\usepackage{xcolor}

% ============================================================
% 自定义数学符号宏
% ============================================================
% 与 config/glossary.yaml 中 symbols 部分保持同步。
% 添加新符号时请同步更新 glossary.yaml。

% \newcommand{\loss}{\mathcal{L}}       % 损失函数
% \newcommand{\params}{\theta}           % 模型参数
% \newcommand{\dataset}{\mathcal{D}}     % 数据集

% ============================================================
% 标记宏 — 用于 Agent 与用户的反馈循环
% ============================================================
\newcommand{\todo}[1]{\textcolor{red}{\textbf{[TODO: #1]}}}
\newcommand{\confirm}[1]{\textcolor{orange}{\textbf{[CONFIRM: #1]}}}
 引入。
% 所有 \usepackage、自定义命令、数学符号宏均在此定义。
% Agent 在 glossary.yaml 中新增符号时，必须同步更新此文件。

% --- 中文支持 ---
\usepackage[UTF8]{ctex}

% --- 数学 ---
\usepackage{amsmath}
\usepackage{amssymb}
\usepackage{amsthm}

% --- 图表 ---
\usepackage{graphicx}
\usepackage{booktabs}
\usepackage{caption}
\usepackage{subcaption}

% --- 引用与链接 ---
\usepackage{hyperref}
\usepackage{cleveref}
\usepackage{natbib}

% --- 其他 ---
\usepackage{algorithm}
\usepackage{algorithmic}
\usepackage{xcolor}

% ============================================================
% 自定义数学符号宏
% ============================================================
% 与 config/glossary.yaml 中 symbols 部分保持同步。
% 添加新符号时请同步更新 glossary.yaml。

% \newcommand{\loss}{\mathcal{L}}       % 损失函数
% \newcommand{\params}{\theta}           % 模型参数
% \newcommand{\dataset}{\mathcal{D}}     % 数据集

% ============================================================
% 标记宏 — 用于 Agent 与用户的反馈循环
% ============================================================
\newcommand{\todo}[1]{\textcolor{red}{\textbf{[TODO: #1]}}}
\newcommand{\confirm}[1]{\textcolor{orange}{\textbf{[CONFIRM: #1]}}}
